\chapter{Convexity and Duality}\label{ch:convexity}

\begin{center}\itshape
In fact the great watershed in optimization isn't between linearity and nonlinearity, but convexity and nonconvexity.\\
\upshape(R.~Tyrrell Rockafellar, 1993)
\end{center}

Some time around 1640, Pierre de Fermat, a magistrate in Toulouse who did mathematics in his letters, sent a challenge to Evangelista Torricelli,
Galileo's successor in Florence: given three points in the plane, find the point whose total distance to the three is smallest. Think of three towns
that want to share a railway station and pay for the track by the mile. Torricelli solved it, and the answer is pure geometry: if every angle of the
triangle is below $120^\circ$, the best point is the one from which the three towns are seen at angles of exactly $120^\circ$ to each other; otherwise it is the
vertex with the obtuse angle. Bonaventura Cavalieri published the $120^\circ$ property in 1647 and Torricelli's pupil Vincenzo Viviani published the
solution in 1659.

Two centuries later, in 1846, a schoolteacher named Fasbender noticed something stranger. Draw any equilateral triangle whose three sides pass through
the three towns. Its height is \emph{never larger} than the total distance from any point to the towns, and the largest such triangle has height
\emph{exactly} equal to the smallest total distance. A minimum problem over points equals a maximum problem over triangles. Harold Kuhn, who dug up
this history in 1974, saw in it the earliest example of a pair of dual nonlinear programs.

This chapter is about that phenomenon. Convex problems come with a mirror problem, the dual, whose values bound the original from the other side,
and under a mild condition the two meet. The meeting point is a certificate: to convince someone that $6.541$ kilometres of track is the least possible,
you do not need to search the plane, you show them a triangle. The tools are separating hyperplanes, the Fenchel conjugate and the Lagrangian, and
the payoffs include the conditions of Karush, Kuhn and Tucker, the water-filling rule for sharing power between noisy channels, the maximum entropy
principle that makes exponential families appear in Chapter~8, and von Neumann's minimax theorem for games. Linear programming duality, proved by
the simplex method in Tome~C2, Chapter~7, reappears as a special case.

\begin{think}
The towns are at $(0,0)$, $(4,0)$ and $(1,3)$. A station at the centre of gravity $(\frac53,1)$ needs $6.611$ units of track. Is the
optimum much better? And what kind of object could \emph{prove} that no station needs less than some number?
\end{think}

\section{Convex sets and separation}\label{sec:separation}

A set $C\subseteq\R^n$ is \emph{convex} if it contains the segment between any two of its points: $x,y\in C$ and $t\in[0,1]$ give $(1-t)x+ty\in C$. Intersections of convex
sets are convex; balls, half-spaces $\{x:a^{\mathsf T}x\le\alpha\}$, polyhedra and the cone of positive semidefinite matrices are convex. The \emph{convex hull} $\operatorname{conv}S$ is
the smallest convex set containing $S$, the set of all convex combinations $\sum_it_ix_i$ of finitely many points of $S$ ($t_i\ge0$, $\sum_it_i=1$).

\begin{lemma}[projection]\label{lem:projection}
Let $C\subseteq\R^n$ be non-empty, closed and convex, and $z\in\R^n$. There is a unique point $p=P_C(z)\in C$ nearest to $z$, and it is characterised by
\[
(z-p)^{\mathsf T}(y-p)\le0\qquad\text{for all }y\in C.
\]
\end{lemma}
\begin{proof}
Fix $y_0\in C$. A nearest point, if any, lies in the compact set $C\cap\{y:\|y-z\|\le\|y_0-z\|\}$, on which the continuous function $\|y-z\|$ attains its minimum: so a nearest point $p$ exists.
For $y\in C$ and $t\in(0,1]$, $p+t(y-p)\in C$, so $\|z-p-t(y-p)\|^2\ge\|z-p\|^2$, that is $-2t(z-p)^{\mathsf T}(y-p)+t^2\|y-p\|^2\ge0$; divide by $t$ and let $t\to0$ to get
the inequality. Conversely, if $p\in C$ satisfies the inequality, then for every $y\in C$,
\[
\|z-y\|^2=\|z-p\|^2+\|p-y\|^2-2(z-p)^{\mathsf T}(y-p)\ge\|z-p\|^2+\|p-y\|^2,
\]
so $p$ is nearest, and any other nearest point $y$ has $\|p-y\|=0$.
\end{proof}

\begin{theorem}[separating a point]\label{thm:pointsep}
If $C\subseteq\R^n$ is non-empty, closed and convex and $z\notin C$, there is $a\in\R^n$ with $a^{\mathsf T}y<a^{\mathsf T}z$ for all $y\in C$; one can take $a=z-P_C(z)$.
\end{theorem}
\begin{proof}
Let $p=P_C(z)$ and $a=z-p\ne0$. By Lemma~\ref{lem:projection}, $a^{\mathsf T}y\le a^{\mathsf T}p=a^{\mathsf T}z-\|a\|^2<a^{\mathsf T}z$.
\end{proof}

To separate two convex sets that may touch, we need to know that the closure of a convex set adds nothing to its interior.

\begin{lemma}\label{lem:intclosure}
If $K\subseteq\R^n$ is convex, then every interior point of $\overline K$ lies in $K$.
\end{lemma}
\begin{proof}
Let $x$ be interior to $\overline K$, with the ball $B(x,r)\subseteq\overline K$. Choose $n+1$ points $w_0,\dots,w_n$ in $B(x,r)$, the vertices of a simplex with $x$ in its interior,
so that $x=\sum_jt_jw_j$ with all barycentric coordinates $t_j>0$. The barycentric coordinates of $x$ with respect to $n+1$ affinely independent points are the solution of
the linear system $\sum_jt_jw_j=x$, $\sum_jt_j=1$, and depend continuously on the points (Cramer's rule). Each $w_j$ is a limit of points of $K$; choosing
$y_j\in K$ close enough to $w_j$, the coordinates of $x$ with respect to $y_0,\dots,y_n$ stay positive, so $x\in\operatorname{conv}\{y_0,\dots,y_n\}\subseteq K$.
\end{proof}

\begin{theorem}[separation of convex sets]\label{thm:separation}
Let $C,D\subseteq\R^n$ be non-empty, convex and disjoint. There are $a\ne0$ and $\alpha\in\R$ with
\[
a^{\mathsf T}x\ge\alpha\ge a^{\mathsf T}y\qquad\text{for all }x\in C,\ y\in D.
\]
\end{theorem}
\begin{proof}
The set $K=C-D=\{x-y:x\in C,y\in D\}$ is convex and does not contain $0$. By Lemma~\ref{lem:intclosure}, $0$ is not interior to $\overline K$, so there are points $z_k\notin\overline K$ with
$z_k\to0$. By Theorem~\ref{thm:pointsep} applied to $\overline K$ and $z_k$, there are unit vectors $a_k$ with $a_k^{\mathsf T}w<a_k^{\mathsf T}z_k$ for all $w\in K$. A subsequence of
$(a_k)$ converges to a unit vector $b$, and in the limit $b^{\mathsf T}w\le0$ on $K$. Put $a=-b$: then $a^{\mathsf T}x\ge a^{\mathsf T}y$ for all $x\in C$, $y\in D$, and
$\alpha=\inf_{x\in C}a^{\mathsf T}x$ is finite and works.
\end{proof}

In infinite dimension the same statement needs one of the sets to be open, and it is the geometric Hahn--Banach theorem (Tome~7, Theorem~1.6). In
$\R^n$ nothing is needed, but separation may not be strict: the disjoint closed convex sets $\{(s,t):s>0,\ st\ge1\}$ and $\{(s,t):t\le0\}$ can be
separated only by the horizontal axis, which touches neither (Exercise~\ref{ex:nostrict}).

\section{Convex functions}\label{sec:convexfunctions}

We allow the value $+\infty$, which encodes constraints: a function $f:\R^n\to(-\infty,+\infty]$ is \emph{convex} if
$f((1-t)x+ty)\le(1-t)f(x)+tf(y)$ for $t\in[0,1]$, equivalently if its \emph{epigraph} $\operatorname{epi}f=\{(x,s):s\ge f(x)\}$ is a convex subset of $\R^{n+1}$. Its
\emph{domain} $\operatorname{dom}f=\{f<\infty\}$ is then convex; $f$ is \emph{proper} if the domain is not empty, and \emph{closed} if the epigraph is closed (that is, $f$ is lower
semicontinuous). The \emph{indicator} $\iota_C$ of a set, $0$ on $C$ and $+\infty$ outside, is convex exactly when $C$ is. Jensen's inequality
$f(\mathbb EX)\le\mathbb Ef(X)$ was proved in Tome~0, Theorem~4.2, and for random variables in Tome~3, Theorem~4.14.

\begin{proposition}[first and second order tests]\label{prop:convextests}
Let $f$ be differentiable on an open convex set $U\subseteq\R^n$.
\begin{enumerate}
\item $f$ is convex on $U$ if and only if $f(y)\ge f(x)+\nabla f(x)^{\mathsf T}(y-x)$ for all $x,y\in U$.
\item If $f$ is twice differentiable, $f$ is convex on $U$ if and only if its Hessian $\nabla^2f(x)$ is positive semidefinite for every $x\in U$.
\end{enumerate}
\end{proposition}
\begin{proof}
Both properties hold on $U$ if and only if they hold on every segment, so it suffices to look at $\phi(t)=f(x+t(y-x))$, $t$ in an open interval containing $[0,1]$, with
$\phi'(t)=\nabla f(x+t(y-x))^{\mathsf T}(y-x)$ and $\phi''(t)=(y-x)^{\mathsf T}\nabla^2f(\cdot)(y-x)$. (1) If $f$ is convex, $\frac{\phi(t)-\phi(0)}t\le\phi(1)-\phi(0)$ for
$t\in(0,1]$, and letting $t\to0$ gives $\phi'(0)\le\phi(1)-\phi(0)$, the inequality. Conversely, if the inequality holds, apply it at $z=(1-t)x+ty$ towards $x$ and towards $y$,
and add the two results with weights $1-t$ and $t$: the gradient terms cancel and $(1-t)f(x)+tf(y)\ge f(z)$. (2) A function of one variable is convex if and only if
its derivative is non-decreasing, which by the mean value theorem holds if and only if $\phi''\ge0$; and $\phi''\ge0$ for all segments is the positive
semidefiniteness of the Hessian.
\end{proof}

So a differentiable convex function lies above all its tangent planes, and a point where the gradient vanishes is a \emph{global} minimum. This is the
watershed of the epigraph: for convex problems, local information certifies global optimality. Functions that are not differentiable, like the distance
$\|x-a\|$ at $x=a$, still have one-sided derivatives.

\begin{lemma}\label{lem:dirderiv}
Let $f:\R^n\to\R$ be convex. For every $x$ and $d$ the \emph{directional derivative} $f'(x;d)=\lim_{t\downarrow0}\frac{f(x+td)-f(x)}t$ exists, and
$f(x+d)\ge f(x)+f'(x;d)$. Hence $x$ minimises $f$ if and only if $f'(x;d)\ge0$ for every direction $d$.
\end{lemma}
\begin{proof}
For $0<s<t$, convexity applied to $x+sd=(1-\frac st)x+\frac st(x+td)$ shows that the quotient $q(t)=\frac{f(x+td)-f(x)}t$ satisfies $q(s)\le q(t)$. Applied to
$x=\frac1{1+t}(x+td)+\frac t{1+t}(x-d)$, it shows $q(t)\ge f(x)-f(x-d)$. A non-decreasing function bounded below has a limit as $t\downarrow0$, and $q(1)\ge\lim q$ is the
inequality. If $f'(x;d)\ge0$ for all $d$, then $f(x+d)\ge f(x)$ for all $d$; the converse is clear.
\end{proof}

\begin{theorem}[Torricelli; Cavalieri, 1647]\label{thm:fermat}
Let $a_1,a_2,a_3$ be three points of the plane, not on a line, and $f(x)=\sum_i\|x-a_i\|$. Then $f$ has a unique minimiser $x_*$.
\begin{enumerate}
\item If the angle of the triangle at $a_1$ is at least $120^\circ$, then $x_*=a_1$.
\item If all angles are below $120^\circ$, $x_*$ is the unique point at which the unit vectors $u_i=(a_i-x_*)/\|a_i-x_*\|$ satisfy $u_1+u_2+u_3=0$, that is, from which each side is
seen under an angle of $120^\circ$.
\end{enumerate}
\end{theorem}
\begin{proof}
$f$ is convex, as a sum of norms of affine functions, and continuous, and $f(x)\ge\|x\|-\sum\|a_i\|$ tends to infinity, so $f$ attains its minimum on a large closed disc. Restricted to a
line, $\|x-a\|$ is strictly convex unless the line passes through $a$ (the graph is then the hyperbola-like curve $t\mapsto\sqrt{t^2+h^2}$ with $h>0$); since no line contains all three
points, $f$ is strictly convex on every line, and two minimisers would give a third, their midpoint, with smaller value. So the minimiser is unique.

At a point $x$ different from the $a_i$, $f$ is differentiable with $\nabla f(x)=-\sum_iu_i$, the $u_i$ being the unit vectors from $x$ towards the $a_i$. It vanishes if and
only if three unit vectors sum to zero, which happens if and only if they make pairwise angles of $120^\circ$: indeed $\|u_1+u_2\|^2=\|u_3\|^2$ reads $2+2u_1^{\mathsf T}u_2=1$.

At $a_1$, put $v_j=(a_j-a_1)/\|a_j-a_1\|$ for $j=2,3$. The term $\|a_1+td-a_1\|=t\|d\|$ and the two differentiable terms give $f'(a_1;d)=\|d\|-(v_2+v_3)^{\mathsf T}d$, which is
$\ge0$ for every $d$ if and only if $\|v_2+v_3\|\le1$, that is $2+2v_2^{\mathsf T}v_3\le1$, an angle at $a_1$ of at least $120^\circ$. By Lemma~\ref{lem:dirderiv}, $a_1$ is the minimiser in case (1).
In case (2) no vertex is a minimiser, so $x_*$ is a point where the gradient vanishes, and it is unique.
\end{proof}

For the towns of the \emph{Pause and think} all angles are below $120^\circ$ (they are $71.6^\circ$, $45^\circ$ and $63.4^\circ$), and the station goes to
$x_*\approx(1.302,1.047)$ with $f(x_*)\approx6.541$: the centre of gravity was only $1\%$ worse. Endre Weiszfeld found in 1937 the iteration
$x\leftarrow\sum_ia_i\|x-a_i\|^{-1}/\sum_i\|x-a_i\|^{-1}$, whose fixed points are exactly the points where the gradient vanishes, and it is how the notebook finds $x_*$
(Kuhn proved in 1973 that it converges, unless it lands on one of the $a_i$).

Now Fasbender's triangle. For any unit vectors, or more generally any $v_1,v_2,v_3$ with $\|v_i\|\le1$ and $v_1+v_2+v_3=0$, and any point $x$, Cauchy--Schwarz gives
\begin{equation}\label{eq:fasbender}
f(x)=\sum_i\|a_i-x\|\ge\sum_iv_i^{\mathsf T}(a_i-x)=\sum_iv_i^{\mathsf T}a_i,
\end{equation}
because $\sum_iv_i^{\mathsf T}x=0$. The left side does not depend on the $v_i$, the right side does not depend on $x$: every admissible triple $(v_i)$ certifies a lower bound for the
whole problem. In case (2), the triple $v_i=u_i$ makes~\eqref{eq:fasbender} an equality at $x=x_*$, so
\[
\min_x f(x)=\max\Bigl\{\sum_iv_i^{\mathsf T}a_i:\ \|v_i\|\le1,\ v_1+v_2+v_3=0\Bigr\}.
\]
When the $v_i$ are unit vectors at $120^\circ$, the lines through $a_i$ perpendicular to $v_i$ bound an equilateral triangle, and $\sum_iv_i^{\mathsf T}a_i$ is its height, by Viviani's
theorem on the distances from an interior point to the sides of an equilateral triangle (Exercise~\ref{ex:viviani}). This is the picture of Figure~\ref{fig:fermat}: the
minimum over stations equals the maximum over circumscribed equilateral triangles. The rest of the chapter turns this trick into a theory.

\begin{figure}[ht]
\centering
\fig{ch03-fermat}
\caption{The Fermat--Torricelli point of the towns $(0,0)$, $(4,0)$, $(1,3)$, with the three roads at $120^\circ$, and Fasbender's dual certificate: the largest equilateral
triangle whose sides pass through the towns, perpendicular to the roads. Its height equals the total length of the roads, $6.541$.}
\label{fig:fermat}
\end{figure}

\section{The Fenchel conjugate}\label{sec:conjugate}

A closed convex set is the intersection of the half-spaces that contain it (Theorem~\ref{thm:pointsep}). The analogue for functions: a closed convex function is the supremum
of the affine functions below it. The affine functions $x\mapsto y^{\mathsf T}x-\beta$ below $f$ with a given slope $y$ are those with $\beta\ge\sup_x(y^{\mathsf T}x-f(x))$, and the best
intercept defines the conjugate.

\begin{definition}\label{def:conjugate}
The \emph{Fenchel conjugate} of $f:\R^n\to(-\infty,+\infty]$ is $f^*(y)=\sup_x\bigl(y^{\mathsf T}x-f(x)\bigr)$, and its \emph{biconjugate} is $f^{**}=(f^*)^*$.
\end{definition}

Whatever $f$ is, $f^*$ is a supremum of affine functions of $y$, hence convex and closed. Directly from the definition comes the \emph{Fenchel--Young inequality}
\[
f(x)+f^*(y)\ge y^{\mathsf T}x,
\]
and so $f^{**}(x)=\sup_y(y^{\mathsf T}x-f^*(y))\le f(x)$. Figure~\ref{fig:conjugate} shows the geometry: $-f^*(y)$ is where the lowest line of slope $y$ touching the graph
crosses the vertical axis. Some conjugates to remember (Exercise~\ref{ex:conjugates}): $\frac12\|x\|^2$ is its own conjugate; the conjugate of $e^x$ is $y\ln y-y$ for $y\ge0$
(with $0\ln0=0$) and $+\infty$ for $y<0$; the conjugate of a norm is the indicator of the unit ball of the dual norm; and the conjugate of the \emph{log-sum-exp} function
$\ln\sum_je^{x_j}$ is the negative entropy $\sum_jy_j\ln y_j$ on the probability simplex, $+\infty$ elsewhere (Exercise~\ref{ex:logsumexp}). The last pair is the heart of the maximum
entropy principle of Section~\ref{sec:kkt}.

\begin{figure}[ht]
\centering
\fig{ch03-conjugate}
\caption{The conjugate read off the graph. For each slope $y$, the line $x\mapsto yx-f^*(y)$ is the highest line of slope $y$ below $f$; it touches the graph where $f'(x)=y$.
Here $f(x)=\frac{x^4}4-x$ and $y=1$, and the supremum of all these lines gives $f$ back: $f^{**}=f$.}
\label{fig:conjugate}
\end{figure}

\begin{theorem}[Fenchel, 1949; Moreau, 1962]\label{thm:fenchelmoreau}
If $f:\R^n\to(-\infty,+\infty]$ is proper, convex and closed, then $f^{**}=f$.
\end{theorem}
\begin{proof}
\emph{Step 1: $f$ has an affine minorant.} Take $x_1\in\operatorname{dom}f$. The point $(x_1,f(x_1)-1)$ is not in the closed convex set $\operatorname{epi}f\subseteq\R^{n+1}$, so by
Theorem~\ref{thm:pointsep} there is $(a,\alpha)\in\R^n\times\R$ with $a^{\mathsf T}x+\alpha s<a^{\mathsf T}x_1+\alpha(f(x_1)-1)$ for all $(x,s)\in\operatorname{epi}f$. Letting $s\to\infty$ shows
$\alpha\le0$, and $(x,s)=(x_1,f(x_1))$ excludes $\alpha=0$. Dividing by $-\alpha>0$: $f(x)\ge b^{\mathsf T}x-\beta$ for all $x$, with $b=-a/\alpha$ and some $\beta\in\R$. In particular
$f^*(b)\le\beta<\infty$.

\emph{Step 2.} We know $f^{**}\le f$. Suppose $f^{**}(x_0)<f(x_0)$ for some $x_0$; then $f^{**}(x_0)$ is finite (it is at least $b^{\mathsf T}x_0-f^*(b)>-\infty$), and the point
$(x_0,f^{**}(x_0))$ lies outside $\operatorname{epi}f$. Separate it strictly: some $(a,\alpha)$ and $c$ satisfy $a^{\mathsf T}x+\alpha s\le c<a^{\mathsf T}x_0+\alpha f^{**}(x_0)$ for all
$(x,s)\in\operatorname{epi}f$, and again $\alpha\le0$.

If $\alpha<0$, scale to $\alpha=-1$: then $a^{\mathsf T}x-f(x)\le c$ for all $x$, that is $f^*(a)\le c$, and $f^{**}(x_0)\ge a^{\mathsf T}x_0-f^*(a)\ge a^{\mathsf T}x_0-c>f^{**}(x_0)$, a
contradiction.

If $\alpha=0$, then $a^{\mathsf T}x\le c<a^{\mathsf T}x_0$ for $x\in\operatorname{dom}f$: the hyperplane is vertical. Tilt the minorant of Step~1 with it: for $\lambda\ge0$ and
$x\in\operatorname{dom}f$, $f(x)\ge b^{\mathsf T}x-\beta+\lambda(a^{\mathsf T}x-c)$, so $f^*(b+\lambda a)\le\beta+\lambda c$ and
\[
f^{**}(x_0)\ge(b+\lambda a)^{\mathsf T}x_0-\beta-\lambda c=b^{\mathsf T}x_0-\beta+\lambda(a^{\mathsf T}x_0-c)\xrightarrow[\lambda\to\infty]{}+\infty,
\]
contradicting the finiteness of $f^{**}(x_0)$.
\end{proof}

For a function that is not convex, $f^{**}$ is its \emph{convex envelope}, the largest closed convex function below it (Exercise~\ref{ex:envelope}): conjugating twice
convexifies, which is how the notebook draws convex envelopes.

\section{Lagrange duality}\label{sec:lagrange}

The general problem of this book's optimisation chapters is
\begin{equation}\label{eq:primalconvex}
\text{minimise }f_0(x)\quad\text{subject to}\quad f_i(x)\le0\ (1\le i\le m),\qquad Ax=b,
\end{equation}
with $f_0,\dots,f_m:\R^n\to\R$ and $A\in\R^{p\times n}$; it is a \emph{convex program} when $f_0,\dots,f_m$ are convex. Let $p^*$ be its optimal value (the infimum over feasible points,
$+\infty$ if there are none). Joseph-Louis Lagrange's idea, in his \emph{M\'ecanique analytique} of 1788, was to price the constraints: the \emph{Lagrangian} is
\[
L(x,\lambda,\nu)=f_0(x)+\sum_{i=1}^m\lambda_if_i(x)+\nu^{\mathsf T}(Ax-b),\qquad\lambda\in\R^m,\ \nu\in\R^p,
\]
and the \emph{dual function} is $g(\lambda,\nu)=\inf_xL(x,\lambda,\nu)$, an infimum of affine functions of $(\lambda,\nu)$, hence concave, possibly $-\infty$.

\begin{proposition}[weak duality]\label{prop:weakdualityc}
For every $\lambda\ge0$ and $\nu$, $g(\lambda,\nu)\le p^*$. Hence the \emph{dual problem}, maximise $g(\lambda,\nu)$ subject to $\lambda\ge0$, has value $d^*\le p^*$.
\end{proposition}
\begin{proof}
If $x$ is feasible and $\lambda\ge0$, then $\lambda_if_i(x)\le0$ and $Ax-b=0$, so $g(\lambda,\nu)\le L(x,\lambda,\nu)\le f_0(x)$. Take the infimum over feasible $x$.
\end{proof}

Weak duality holds for every problem, convex or not, and it is often all one needs: a dual feasible point is a certificate of a lower bound, as Fasbender's
triangle was. The difference $p^*-d^*$ is the \emph{duality gap}. For convex programs it vanishes under a condition that is usually easy to check.

\begin{theorem}[Slater, 1950]\label{thm:slater}
Let~\eqref{eq:primalconvex} be a convex program with $p^*$ finite, and suppose there is a \emph{strictly feasible} point: $\bar x$ with $f_i(\bar x)<0$ for all $i$ and $A\bar x=b$. Then
$d^*=p^*$, and the dual optimum is attained: there are $\lambda^*\ge0$ and $\nu^*$ with $g(\lambda^*,\nu^*)=p^*$.
\end{theorem}
\begin{proof}
Since $\bar x$ is feasible, the system $Ax=b$ is consistent, and removing equations that are linear combinations of others changes neither the feasible set nor
$p^*$; a dual solution of the reduced problem extends by zero multipliers. So we may assume that $A$ has full row rank. Consider the convex sets of $\R^m\times\R^p\times\R$
\[
\mathcal A=\{(u,v,t):\ \exists x\ \ f_i(x)\le u_i\ (\forall i),\ Ax-b=v,\ f_0(x)\le t\},\qquad\mathcal B=\{(0,0,s):s<p^*\}.
\]
They are disjoint, since a point of $\mathcal A$ with $u\le0$, $v=0$ comes from a feasible $x$ with $f_0(x)\le t$, so $t\ge p^*$. By Theorem~\ref{thm:separation} there are
$(\lambda,\nu,\mu)\ne0$ and $\alpha$ with $\lambda^{\mathsf T}u+\nu^{\mathsf T}v+\mu t\ge\alpha$ on $\mathcal A$ and $\mu s\le\alpha$ for all $s<p^*$. As $\mathcal A$ is unchanged when $u$ or $t$ increase,
$\lambda\ge0$ and $\mu\ge0$; and $\mu p^*\le\alpha$. Taking $u_i=f_i(x)$, $v=Ax-b$, $t=f_0(x)$:
\begin{equation}\label{eq:slaterkey}
\mu f_0(x)+\sum_i\lambda_if_i(x)+\nu^{\mathsf T}(Ax-b)\ge\mu p^*\qquad\text{for all }x.
\end{equation}
If $\mu>0$, divide by $\mu$: $L(x,\lambda/\mu,\nu/\mu)\ge p^*$ for all $x$, so $g(\lambda/\mu,\nu/\mu)\ge p^*$, and weak duality gives equality. If $\mu=0$, \eqref{eq:slaterkey} at
$x=\bar x$ gives $\sum_i\lambda_if_i(\bar x)\ge0$ with every $f_i(\bar x)<0$ and $\lambda\ge0$, so $\lambda=0$; then $\nu\ne0$ and $\nu^{\mathsf T}(Ax-b)\ge0$ for all $x$. An affine function
bounded below on $\R^n$ is constant, so $A^{\mathsf T}\nu=0$, which for $A$ of full row rank forces $\nu=0$: a contradiction.
\end{proof}

Without a strictly feasible point the theorem can fail, and Exercise~\ref{ex:slaterfail} shows two ways. For linear programs, however, no Slater condition is needed: a
feasible linear program with finite value always has zero duality gap (Tome~C2, Theorem~7.7, proved there by the simplex method). Let us see the Lagrangian
produce the linear programming dual. For the program, minimise $c^{\mathsf T}x$ subject to $Ax=b$ and $-x\le0$,
\[
g(\mu,\nu)=\inf_x\bigl((c-\mu+A^{\mathsf T}\nu)^{\mathsf T}x-b^{\mathsf T}\nu\bigr)=\begin{cases}-b^{\mathsf T}\nu&\text{if }c-\mu+A^{\mathsf T}\nu=0,\\-\infty&\text{otherwise.}\end{cases}
\]
With $y=-\nu$ and $\mu=c-A^{\mathsf T}y\ge0$, the dual is: maximise $b^{\mathsf T}y$ subject to $A^{\mathsf T}y\le c$, the dual of Tome~C2, Chapter~7, in another standard form.

\section{The conditions of Karush, Kuhn and Tucker}\label{sec:kkt}

When the functions are differentiable, zero duality gap becomes a system of equations and inequalities. William Karush wrote it down in his 1939 master's
thesis at the University of Chicago, which nobody read; Harold Kuhn and Albert Tucker found it again in 1951, and the conditions now carry all three names.

\begin{theorem}[Karush, 1939; Kuhn and Tucker, 1951]\label{thm:kkt}
Let~\eqref{eq:primalconvex} be a convex program with $f_0,\dots,f_m$ differentiable. Consider the \emph{KKT conditions} for $x^*\in\R^n$, $\lambda^*\in\R^m$, $\nu^*\in\R^p$:
\begin{enumerate}
\item[(a)] primal feasibility: $f_i(x^*)\le0$ for all $i$, and $Ax^*=b$;
\item[(b)] dual feasibility: $\lambda^*\ge0$;
\item[(c)] complementary slackness: $\lambda_i^*f_i(x^*)=0$ for all $i$;
\item[(d)] stationarity: $\nabla f_0(x^*)+\sum_i\lambda_i^*\nabla f_i(x^*)+A^{\mathsf T}\nu^*=0$.
\end{enumerate}
If they hold, $x^*$ is optimal, $(\lambda^*,\nu^*)$ is dual optimal and $p^*=d^*$. Conversely, if $x^*$ is optimal and Slater's condition holds, there are $\lambda^*,\nu^*$
satisfying the KKT conditions.
\end{theorem}
\begin{proof}
Suppose (a)--(d). Since $\lambda^*\ge0$, $x\mapsto L(x,\lambda^*,\nu^*)$ is convex and differentiable, and by (d) its gradient vanishes at $x^*$; by Proposition~\ref{prop:convextests}(1)
it is minimised there. So $g(\lambda^*,\nu^*)=L(x^*,\lambda^*,\nu^*)=f_0(x^*)$ by (a) and (c). Weak duality, $g(\lambda^*,\nu^*)\le p^*\le f_0(x^*)$, turns into equalities: $x^*$ is primal
optimal, $(\lambda^*,\nu^*)$ dual optimal, and the gap is zero.

Conversely, let $x^*$ be optimal and $(\lambda^*,\nu^*)$ a dual optimum given by Theorem~\ref{thm:slater}. Then
\[
f_0(x^*)=g(\lambda^*,\nu^*)\le L(x^*,\lambda^*,\nu^*)=f_0(x^*)+\sum_i\lambda_i^*f_i(x^*)\le f_0(x^*),
\]
so both inequalities are equalities. The second says $\sum_i\lambda_i^*f_i(x^*)=0$, a sum of non-positive terms, hence (c). The first says that $x^*$ minimises the differentiable
function $L(\cdot,\lambda^*,\nu^*)$ on $\R^n$, hence (d).
\end{proof}

The multiplier $\lambda_i^*$ is the price of the $i$th constraint: complementary slackness says that a constraint that is not tight costs nothing. Two examples show the
conditions at work; each time they do more than verify a solution, they \emph{find} it.

\paragraph{Water-filling.} A transmitter has total power $P>0$ to share between $n$ parallel channels with noise levels $\alpha_1,\dots,\alpha_n>0$. Putting power
$x_i$ in channel $i$ yields a rate proportional to $\ln(\alpha_i+x_i)-\ln\alpha_i$ (for Gaussian channels this is Shannon's capacity formula; see Cover and Thomas, Section~9.4,
and Tome~C4, Chapter~3, for channel capacity). The transmitter wants to
\[
\text{maximise}\ \sum_{i=1}^n\ln(\alpha_i+x_i)\quad\text{subject to}\quad x\ge0,\ \sum_ix_i=P.
\]

\begin{proposition}[water-filling]\label{prop:waterfilling}
There is a unique level $w>0$ with $\sum_i\max(0,w-\alpha_i)=P$, and the unique optimal allocation is $x_i^*=\max(0,w-\alpha_i)$.
\end{proposition}
\begin{proof}
The function $w\mapsto\sum_i\max(0,w-\alpha_i)$ is continuous, zero for $w\le\min\alpha_i$, strictly increasing after that and unbounded, so the level $w$ exists and is unique. Write the
problem as minimising the convex function $-\sum_i\ln(\alpha_i+x_i)$ (with value $+\infty$ when some $\alpha_i+x_i\le0$; it is differentiable on the open set where all are positive, which
contains the feasible set) under $-x_i\le0$ and $\sum_ix_i=P$. We check the KKT conditions for $x^*$, $\nu^*=1/w$ and $\lambda_i^*=1/w-1/(\alpha_i+x_i^*)$. Stationarity reads
$-\frac1{\alpha_i+x_i^*}-\lambda_i^*+\nu^*=0$, true by the choice of $\lambda^*$. If $x_i^*>0$, then $\alpha_i+x_i^*=w$ and $\lambda_i^*=0$; if $x_i^*=0$, then $\alpha_i\ge w$ and
$\lambda_i^*=\frac1w-\frac1{\alpha_i}\ge0$. So $\lambda^*\ge0$ and $\lambda_i^*x_i^*=0$, and Theorem~\ref{thm:kkt} shows that $x^*$ is optimal. The objective is strictly concave, so the optimum is
unique.
\end{proof}

Picture the noise levels as the uneven floor of a tank and pour in a volume $P$ of water: the water surface is flat at height $w$, and the depth above channel $i$
is its power (Figure~\ref{fig:waterfilling}). Channels that are too noisy stay dry.

\begin{figure}[ht]
\centering
\fig{ch03-waterfilling}
\caption{Water-filling for noise levels $\alpha=(0.5,1.6,0.9,2.5,0.3,1.2,3.0,0.8)$ and power $P=5$: the water level is $w=10.3/6\approx1.717$, and the two noisiest channels receive nothing.}
\label{fig:waterfilling}
\end{figure}

\paragraph{Maximum entropy.} In 1957 Edwin Jaynes proposed a rule for choosing a probability distribution when all one knows are some averages: among all
distributions with those averages, take the one with the largest entropy $H(p)=-\sum_\omega p_\omega\ln p_\omega$, the least committal one. He showed that it reproduces the Gibbs
distributions of statistical mechanics, and his favourite toy was a die whose average throw is known to be $4.5$ instead of $3.5$. Let $\Omega$ be a finite set,
$\phi:\Omega\to\R^K$ a vector of observables and $c\in\R^K$ the known averages. For $\theta\in\R^K$ put
\[
Z(\theta)=\sum_{\omega\in\Omega}e^{\theta^{\mathsf T}\phi(\omega)},\qquad p_\theta(\omega)=\frac{e^{\theta^{\mathsf T}\phi(\omega)}}{Z(\theta)}.
\]
The distributions $p_\theta$ form an \emph{exponential family}, and $\ln Z$ is a log-sum-exp function, convex in $\theta$.

\begin{theorem}[maximum entropy; Jaynes, 1957]\label{thm:maxent}
\begin{enumerate}
\item If $p$ is a probability on $\Omega$ with $\sum_\omega p_\omega\phi(\omega)=c$, then $H(p)\le\ln Z(\theta)-\theta^{\mathsf T}c$ for every $\theta\in\R^K$.
\item If some $p_\theta$ has averages $c$, it is the unique distribution of maximum entropy with averages $c$, and $H(p_\theta)=\ln Z(\theta)-\theta^{\mathsf T}c=\min_{\theta'}\bigl(\ln Z(\theta')-\theta'^{\mathsf T}c\bigr)$.
\end{enumerate}
\end{theorem}
\begin{proof}
For probabilities $p$ and $q$ with $q>0$, \emph{Gibbs' inequality} $\sum_\omega p_\omega\ln\frac{p_\omega}{q_\omega}\ge0$ holds, with equality only if $p=q$: by $\ln u\le u-1$, with equality only at
$u=1$, $-\sum_\omega p_\omega\ln\frac{q_\omega}{p_\omega}\ge-\sum_\omega p_\omega(\frac{q_\omega}{p_\omega}-1)\ge0$, the sums running over $p_\omega>0$. With $q=p_\theta$ and $\ln p_\theta(\omega)=\theta^{\mathsf T}\phi(\omega)-\ln Z(\theta)$,
\[
H(p)\le-\sum_\omega p_\omega\ln p_\theta(\omega)=\ln Z(\theta)-\theta^{\mathsf T}\sum_\omega p_\omega\phi(\omega)=\ln Z(\theta)-\theta^{\mathsf T}c,
\]
which is (1), with equality if and only if $p=p_\theta$. If $p_\theta$ itself has averages $c$, the same computation gives $H(p_\theta)=\ln Z(\theta)-\theta^{\mathsf T}c$, which by (1) is at most
$\ln Z(\theta')-\theta'^{\mathsf T}c$ for all $\theta'$ and at least $H(p)$ for all admissible $p$, with equality only for $p=p_\theta$.
\end{proof}

Part (1) is weak duality and part (2) is zero gap: the dual of entropy maximisation is the minimisation of the convex function $\ln Z(\theta)-\theta^{\mathsf T}c$, exactly the
Lagrange dual (Exercise~\ref{ex:logsumexp} computes it). Its gradient vanishes where $\nabla\ln Z(\theta)=\sum_\omega p_\theta(\omega)\phi(\omega)=c$, the moment equation, which the notebook
solves by Newton's method. For Jaynes' die, $\phi(\omega)=\omega$ on $\{1,\dots,6\}$ and $c=4.5$: the answer is $p_\omega\propto e^{\theta\omega}$ with $\theta\approx0.371$, the probabilities
$0.054$, $0.079$, $0.114$, $0.165$, $0.240$, $0.347$. The same theorem with the mean and the variance as constraints is why the Gaussian appears everywhere
(Exercise~\ref{ex:maxentgauss}), and Chapter~8 will estimate the parameter $\theta$ of an exponential family from data.

\section{Games and the minimax theorem}\label{sec:minimax}

Two players choose simultaneously: Rose a row $i\in\{1,\dots,m\}$, Colin a column $j\in\{1,\dots,n\}$, and Rose pays Colin $M_{ij}$. Mixed strategies are probability vectors
$x\in\Delta_m$ and $y\in\Delta_n$ (the simplices), and the expected payment is $x^{\mathsf T}My$. If Rose announces $x$, Colin answers with the best column, and she pays
$\max_yx^{\mathsf T}My$; Rose minimises this. If Colin announces first, he gets $\min_xx^{\mathsf T}My$. Announcing first can only hurt:
\[
\max_{y\in\Delta_n}\min_{x\in\Delta_m}x^{\mathsf T}My\le\min_{x\in\Delta_m}\max_{y\in\Delta_n}x^{\mathsf T}My,
\]
because $\min_{x'}x'^{\mathsf T}My\le x^{\mathsf T}My\le\max_{y'}x^{\mathsf T}My'$ for any $x,y$. \'Emile Borel asked in 1921 whether equality always holds, and believed it would fail
for large games; John von Neumann proved in 1928 that it never fails. The order of announcement does not matter.

\begin{theorem}[von Neumann, 1928]\label{thm:minimax}
For every real $m\times n$ matrix $M$, $\displaystyle\max_{y\in\Delta_n}\min_{x\in\Delta_m}x^{\mathsf T}My=\min_{x\in\Delta_m}\max_{y\in\Delta_n}x^{\mathsf T}My$, and both extrema are attained.
\end{theorem}
\begin{proof}
A linear function on a simplex reaches its maximum at a vertex, so $\max_yx^{\mathsf T}My=\max_j(M^{\mathsf T}x)_j$. Rose's problem is therefore the convex program
\[
\text{minimise }t\quad\text{over }(x,t)\in\R^m\times\R\quad\text{subject to}\quad(M^{\mathsf T}x)_j-t\le0\ (j\le n),\quad-x_i\le0\ (i\le m),\quad\textstyle\sum_ix_i=1.
\]
Its value $p^*$ is finite (it is a continuous function minimised on the compact simplex) and attained, and $x_i=\frac1m$, $t=1+\max_j(M^{\mathsf T}x)_j$ is strictly feasible. By Theorem~\ref{thm:slater},
the dual has value $p^*$ and is attained. With multipliers $y\in\R^n$, $\mu\in\R^m$ and $\nu\in\R$,
\[
L=t\Bigl(1-\sum_jy_j\Bigr)+x^{\mathsf T}(My-\mu+\nu\mathbf1)-\nu,
\]
whose infimum over $(x,t)$ is $-\nu$ if $\sum_jy_j=1$ and $My-\mu+\nu\mathbf1=0$, and $-\infty$ otherwise. Given $y\ge0$ with $\sum_jy_j=1$, the best choice is the smallest $\nu$ with
$\mu=My+\nu\mathbf1\ge0$, namely $\nu=-\min_i(My)_i$. So the dual problem is to maximise $\min_i(My)_i=\min_xx^{\mathsf T}My$ over $y\in\Delta_n$, and its value is $p^*$.
\end{proof}

So every finite zero-sum game has a \emph{value} and optimal mixed strategies for both players. Tome~C2 met this theorem twice, as a consequence of linear programming
duality (Exercise~7.7 there) and through two Hedge players learning against each other (Chapter~12 there); Chapter~\ref{ch:convexity}'s proof is the third, and
Chapter~7 of this book will turn the second into a general method. Von Neumann's own proof was topological; the one through separation is due to Jean Ville
(1938), and the generalisation to convex-concave functions on compact convex sets is Sion's theorem (1958).

\begin{beauty}
Duality is a theorem about seeing a convex set from outside. A closed convex set is the intersection of its supporting half-spaces, a closed convex function the
supremum of its affine minorants, and an optimisation problem is bounded below by every price system on its constraints. That these outside descriptions
are complete, that nothing hides between the supporting hyperplanes, is the separation theorem, and every certificate in this chapter, Fasbender's triangle,
the water level, Jaynes' temperature and von Neumann's mixed strategy, is a supporting hyperplane in disguise.
\end{beauty}

\section{Exercises}

\begin{exercise}[$\star$]\label{ex:nonexpansive}
Show that the projection onto a non-empty closed convex set is non-expansive: $\|P_C(z)-P_C(z')\|\le\|z-z'\|$.
\end{exercise}

\begin{exercise}[$\star$]\label{ex:nostrict}
Show that the closed convex sets $C=\{(s,t):s>0,\ st\ge1\}$ and $D=\{(s,t):t\le0\}$ are disjoint, that the only hyperplane separating them is the horizontal axis, and that
no hyperplane separates them strictly. Why does Theorem~\ref{thm:pointsep} not apply?
\end{exercise}

\begin{exercise}[$\star\star$]\label{ex:viviani}
(Viviani.) Show that the sum of the distances from a point inside an equilateral triangle to its three sides equals the height. Deduce the geometric meaning
of~\eqref{eq:fasbender}: if $v_1,v_2,v_3$ are unit vectors at $120^\circ$ and the lines through $a_i$ perpendicular to $v_i$ bound an equilateral triangle containing the $a_i$, then
$\sum_iv_i^{\mathsf T}a_i$ is its height.
\end{exercise}

\begin{exercise}[$\star\star$]\label{ex:weiszfeld}
(Weighted Fermat problem.) For points $a_1,\dots,a_k$ and weights $w_i>0$, show that a point $x$ different from the $a_i$ minimises $\sum_iw_i\|x-a_i\|$ if and only if it is a fixed
point of Weiszfeld's map $T(x)=\sum_iw_ia_i\|x-a_i\|^{-1}/\sum_iw_i\|x-a_i\|^{-1}$, and that $a_1$ is a minimiser if and only if $\|\sum_{i\ge2}w_i\frac{a_i-a_1}{\|a_i-a_1\|}\|\le w_1$.
\end{exercise}

\begin{exercise}[$\star$]\label{ex:conjugates}
Compute the conjugates of: (a) $\frac12x^{\mathsf T}Qx$ with $Q$ symmetric positive definite; (b) $e^x$ on $\R$; (c) a norm $\|\cdot\|$ on $\R^n$, in terms of the dual norm
$\|y\|_*=\max_{\|x\|\le1}y^{\mathsf T}x$; (d) $\iota_C$ for a convex set $C$ (the \emph{support function}).
\end{exercise}

\begin{exercise}[$\star\star$]\label{ex:logsumexp}
Let $\operatorname{lse}(x)=\ln\sum_{j=1}^ne^{x_j}$. Show that its Hessian is $\operatorname{diag}(p)-pp^{\mathsf T}$ with $p_j=e^{x_j}/\sum_ke^{x_k}$, hence positive semidefinite. Show that
$\operatorname{lse}^*(y)=\sum_jy_j\ln y_j$ if $y\in\Delta_n$ and $+\infty$ otherwise. Deduce that the Lagrange dual function of the maximum entropy problem of Theorem~\ref{thm:maxent} is
$\ln Z(\theta)-\theta^{\mathsf T}c$.
\end{exercise}

\begin{exercise}[$\star\star$]\label{ex:envelope}
Show that for every $f:\R^n\to(-\infty,+\infty]$ having an affine minorant, $f^{**}$ is the largest closed convex function below $f$. Compute $f^{**}$ for
$f(x)=\min\bigl((x-1)^2,(x+1)^2\bigr)$ on $\R$.
\end{exercise}

\begin{exercise}[$\star\star$]\label{ex:slaterfail}
(a) For the problem, minimise $x$ subject to $x^2\le0$, show that $p^*=d^*=0$ but the dual optimum is not attained. (b) For the problem in $\R^2$, minimise $e^{-x_1}$ subject to
$f_1(x)\le0$ where $f_1(x)=x_1^2/x_2$ for $x_2>0$ and $+\infty$ otherwise (a convex function), show that $p^*=1$ and $d^*=0$.
\end{exercise}

\begin{exercise}[$\star$]\label{ex:leastnorm}
Let $A\in\R^{p\times n}$ have full row rank. Use the KKT conditions to show that the solution of $Ax=b$ of least norm is $x^*=A^{\mathsf T}(AA^{\mathsf T})^{-1}b$, and identify the multiplier.
\end{exercise}

\begin{exercise}[$\star$]\label{ex:waterlevel}
In Proposition~\ref{prop:waterfilling}, show that the optimal value is a concave, increasing function of $P$ whose derivative is $1/w$, the multiplier $\nu^*$. (The multiplier is
the marginal value of power.)
\end{exercise}

\begin{exercise}[$\star\star$]\label{ex:maxentgauss}
Let $p$ be a probability density on $\R$ with mean $0$ and variance $\sigma^2$, and $h(p)=-\int p\ln p$ its differential entropy (assumed finite). Show that
$h(p)\le\frac12\ln(2\pi e\sigma^2)$, with equality only for the Gaussian density. (Imitate the proof of Theorem~\ref{thm:maxent}.)
\end{exercise}

\begin{exercise}[$\star$]\label{ex:saddle}
Let $\phi:X\times Y\to\R$ be any function. Show that $\max_y\min_x\phi=\min_x\max_y\phi$, with all extrema attained, holds if and only if $\phi$ has a saddle point: $(x^*,y^*)$ with
$\phi(x^*,y)\le\phi(x^*,y^*)\le\phi(x,y^*)$ for all $x,y$.
\end{exercise}

\begin{exercise}[$\star$]\label{ex:rps}
Find the value and all optimal strategies of rock--paper--scissors (payment $1$ for a win, $0$ for a tie), and of the game where Rose and Colin each show one or two fingers
and Rose pays Colin the total number of fingers if it is even and receives it if it is odd.
\end{exercise}

\section{Hints}

\begin{hint}{ex:nonexpansive}
Write the characterisation of Lemma~\ref{lem:projection} for $z$ with $y=P_C(z')$ and for $z'$ with $y=P_C(z)$, and add.
\end{hint}
\begin{hint}{ex:nostrict}
The points $(s,1/s)$ approach the axis. For Theorem~\ref{thm:pointsep}, a point is compact; here neither set is.
\end{hint}
\begin{hint}{ex:viviani}
Cut the triangle into three triangles with apex at the point.
\end{hint}
\begin{hint}{ex:weiszfeld}
The gradient condition $\sum_iw_i\frac{x-a_i}{\|x-a_i\|}=0$ can be solved for the $x$ in the numerators. At $a_1$ compute the directional derivative as in Theorem~\ref{thm:fermat}.
\end{hint}
\begin{hint}{ex:conjugates}
For (a), maximise a concave quadratic. For (c), if $\|y\|_*>1$ take $x$ along a direction where $y^{\mathsf T}x>\|x\|$ and scale.
\end{hint}
\begin{hint}{ex:logsumexp}
For the Hessian, $v^{\mathsf T}(\operatorname{diag}(p)-pp^{\mathsf T})v$ is a variance. For the conjugate, set the gradient $p(x)$ equal to $y$; outside the simplex let $x$ run to infinity along a suitable
direction.
\end{hint}
\begin{hint}{ex:envelope}
$f^{**}$ is the supremum of the affine minorants of $f$. Apply Theorem~\ref{thm:fenchelmoreau} to a closed convex $h\le f$.
\end{hint}
\begin{hint}{ex:slaterfail}
(a) $g(\lambda)=-1/(4\lambda)$. (b) The only feasible points have $x_1=0$; and $g(\lambda)=\inf_x(e^{-x_1}+\lambda x_1^2/x_2)$: let $x_2$ grow faster than $x_1^2$.
\end{hint}
\begin{hint}{ex:leastnorm}
Minimise $\frac12\|x\|^2$ subject to $Ax=b$; stationarity is $x+A^{\mathsf T}\nu=0$.
\end{hint}
\begin{hint}{ex:waterlevel}
The level $w(P)$ is increasing and continuous; differentiate $\sum_{i:\alpha_i<w}\ln w+\sum_{i:\alpha_i\ge w}\ln\alpha_i$.
\end{hint}
\begin{hint}{ex:maxentgauss}
Use $\int p\ln(p/q)\ge0$ with $q$ the Gaussian density, and $\ln q(x)=-\frac12\ln(2\pi\sigma^2)-\frac{x^2}{2\sigma^2}$.
\end{hint}
\begin{hint}{ex:saddle}
Take $x^*$ attaining the min--max and $y^*$ attaining the max--min.
\end{hint}
\begin{hint}{ex:rps}
For the finger game, guess an optimal strategy for Rose that makes Colin indifferent between his two columns.
\end{hint}

\section{Solutions}

\begin{solution}{ex:nonexpansive}
\idea The characterisation says that $z-P_C(z)$ points away from $C$; two such inequalities, added, compare the projections with the points.
\writeup Let $p=P_C(z)$ and $p'=P_C(z')$. Lemma~\ref{lem:projection} gives $(z-p)^{\mathsf T}(p'-p)\le0$ and $(z'-p')^{\mathsf T}(p-p')\le0$. Adding,
$\bigl((z-z')-(p-p')\bigr)^{\mathsf T}(p'-p)\le0$, that is $\|p-p'\|^2\le(z-z')^{\mathsf T}(p-p')\le\|z-z'\|\,\|p-p'\|$ by Cauchy--Schwarz. Divide by $\|p-p'\|$ (if it is not zero).
\end{solution}

\begin{solution}{ex:nostrict}
\idea The hyperbola is asymptotic to the axis, so any separating line must be squeezed onto it.
\writeup Points of $C$ have $t\ge1/s>0$ and points of $D$ have $t\le0$: disjoint. Both sets are closed and convex ($C$ is the epigraph of the convex function $1/s$ on $s>0$, together with
its domain). A separating line $\{a_1s+a_2t=\alpha\}$ with $a_1s+a_2t\ge\alpha$ on $C$ and $\le\alpha$ on $D$: since $D$ contains $(s,0)$ for all $s\in\R$, $a_1=0$; then $a_2\ne0$, and $D$
containing $(0,t)$ for all $t\le0$ forces $a_2>0$, so the line is $t=\alpha/a_2$ with $\alpha\ge0$ (from $D$) and $\alpha/a_2\le\inf_Ct=0$ (from $C$): it is the axis $t=0$. It touches the closure
of the gap: $\operatorname{dist}(C,D)=\inf_{s>0}1/s=0$, so no strict separation is possible. Theorem~\ref{thm:pointsep} separates a closed convex set from a \emph{point}; the
proof of strict separation between two sets would need $C-D$ to be closed, and here $C-D=\{t>0\}\cup\dots$ is not: $0$ is in its closure but not in it.
\end{solution}

\begin{solution}{ex:viviani}
\idea Area, computed two ways.
\writeup Let the triangle have side $s$ and height $h$, and let $d_1,d_2,d_3$ be the distances from an interior point $x$ to the sides. Joining $x$ to the vertices cuts the triangle into three
triangles with bases the sides and heights the $d_i$: the area is $\frac12s(d_1+d_2+d_3)=\frac12sh$, so $d_1+d_2+d_3=h$. Now let $v_i$ be unit vectors at $120^\circ$ and $\ell_i$ the line
through $a_i$ perpendicular to $v_i$, and suppose they bound an equilateral triangle $T$ with the $v_i$ pointing outwards. For $x$ inside $T$, the distance from $x$ to $\ell_i$ is
$v_i^{\mathsf T}(a_i-x)$. Summing and using $\sum_iv_i=0$, the height of $T$ is $\sum_iv_i^{\mathsf T}a_i$: the right side of~\eqref{eq:fasbender} is the height of a circumscribed
equilateral triangle, and the duality of Section~\ref{sec:convexfunctions} says that the largest one has height $\min f$.
\end{solution}

\begin{solution}{ex:weiszfeld}
\idea Away from the $a_i$, the function is differentiable, and the zero-gradient equation, read as a fixed point equation, is Weiszfeld's map.
\writeup Let $F(x)=\sum_iw_i\|x-a_i\|$, convex. At $x\ne a_i$, $\nabla F(x)=\sum_iw_i(x-a_i)/\|x-a_i\|$, which vanishes if and only if $x\sum_iw_i\|x-a_i\|^{-1}=\sum_iw_ia_i\|x-a_i\|^{-1}$, that
is $x=T(x)$; and for a convex differentiable function a zero gradient is equivalent to a minimum (Proposition~\ref{prop:convextests}). At $a_1$, with
$v_i=(a_i-a_1)/\|a_i-a_1\|$, $F'(a_1;d)=w_1\|d\|-\sum_{i\ge2}w_iv_i^{\mathsf T}d$, and $\min_{\|d\|=1}F'(a_1;d)=w_1-\|\sum_{i\ge2}w_iv_i\|$ (take $d$ along $\sum_{i\ge2}w_iv_i$). By
Lemma~\ref{lem:dirderiv}, $a_1$ is a minimiser if and only if this is $\ge0$.
\end{solution}

\begin{solution}{ex:conjugates}
\idea Each supremum is attained where the derivative of $y^{\mathsf T}x-f(x)$ vanishes, or not at all.
\writeup (a) $y^{\mathsf T}x-\frac12x^{\mathsf T}Qx$ is a concave quadratic maximised at $x=Q^{-1}y$, with value $\frac12y^{\mathsf T}Q^{-1}y$; for $Q=I$, $\frac12\|x\|^2$ is its own conjugate. (b) $yx-e^x$: for
$y>0$ the maximum is at $x=\ln y$, value $y\ln y-y$; for $y=0$ the supremum is $0$; for $y<0$, $x\to-\infty$ gives $+\infty$. (c) If $\|y\|_*\le1$, then $y^{\mathsf T}x\le\|y\|_*\|x\|\le\|x\|$ and the
supremum of $y^{\mathsf T}x-\|x\|$ is $0$ (at $x=0$). If $\|y\|_*>1$, pick $x_0$ with $\|x_0\|\le1$ and $y^{\mathsf T}x_0>1$; then $y^{\mathsf T}(tx_0)-\|tx_0\|\ge t(y^{\mathsf T}x_0-1)\to\infty$. So $\|\cdot\|^*$ is the
indicator of the dual unit ball. (d) $\iota_C^*(y)=\sup_{x\in C}y^{\mathsf T}x$, the support function of $C$; by Theorem~\ref{thm:fenchelmoreau} it determines $\iota_C$, hence $C$, when $C$ is closed.
\end{solution}

\begin{solution}{ex:logsumexp}
\idea The Hessian of log-sum-exp is a covariance matrix; the gradient is a probability vector, and inverting the gradient map gives the conjugate.
\writeup With $S=\sum_ke^{x_k}$, $\partial_j\operatorname{lse}=e^{x_j}/S=p_j$ and $\partial_{jk}\operatorname{lse}=p_j\delta_{jk}-p_jp_k$. For $v\in\R^n$,
$v^{\mathsf T}(\operatorname{diag}(p)-pp^{\mathsf T})v=\sum_jp_jv_j^2-(\sum_jp_jv_j)^2$, the variance of $v_J$ when $J$ has law $p$: non-negative. For the conjugate: if some $y_j<0$, let $x=-te_j$,
$t\to\infty$: $y^{\mathsf T}x-\operatorname{lse}(x)\ge-ty_j-\ln n\to\infty$. If $y\ge0$ but $\sum_jy_j\ne1$, let $x=t\mathbf1$: $y^{\mathsf T}x-\operatorname{lse}(x)=t(\sum_jy_j-1)-\ln n$, unbounded as
$t\to\pm\infty$. If $y\in\Delta_n$ with all $y_j>0$, the concave function $y^{\mathsf T}x-\operatorname{lse}(x)$ has zero gradient where $p(x)=y$, for instance at $x_j=\ln y_j$, with value
$\sum_jy_j\ln y_j-\ln1$. If some $y_j=0$, the same value is the limit of the values at $y_\varepsilon=(1-\varepsilon)y+\varepsilon\mathbf1/n$, and the supremum of the closed convex function $\operatorname{lse}^*$
along this segment is attained at the end by lower semicontinuity and convexity; so $\operatorname{lse}^*(y)=\sum_jy_j\ln y_j$ on $\Delta_n$.

For maximum entropy, maximise $H(p)$ subject to $\sum_\omega p_\omega\phi(\omega)=c$ over $p\in\Delta_\Omega$, keeping the simplex as the domain. With multiplier $\theta$ for the moment constraint,
the dual function is $\sup_{p\in\Delta}\bigl(H(p)+\theta^{\mathsf T}(\sum_\omega p_\omega\phi(\omega)-c)\bigr)=\sup_p\bigl(\sum_\omega p_\omega\theta^{\mathsf T}\phi(\omega)-\sum_\omega p_\omega\ln p_\omega\bigr)-\theta^{\mathsf T}c$, and the
supremum is $(\text{negative entropy})^*$ at the vector $(\theta^{\mathsf T}\phi(\omega))_\omega$, which by Theorem~\ref{thm:fenchelmoreau} applied to $\operatorname{lse}$ is $\operatorname{lse}(\theta^{\mathsf T}\phi)=\ln Z(\theta)$.
\end{solution}

\begin{solution}{ex:envelope}
\idea The biconjugate is built only from affine functions below $f$, and a closed convex function is the supremum of such functions.
\writeup By definition $f^{**}(x)=\sup_y(y^{\mathsf T}x-f^*(y))$, and $x\mapsto y^{\mathsf T}x-\beta$ is below $f$ if and only if $\beta\ge f^*(y)$: so $f^{**}$ is the supremum of all affine minorants of $f$.
It is closed, convex, below $f$, and proper (it is at least the given affine minorant). Let $h\le f$ be closed and convex with values in $(-\infty,+\infty]$. If $h\equiv+\infty$, then $f\equiv+\infty$ and there is nothing to prove. Otherwise $h$ is proper, and
Theorem~\ref{thm:fenchelmoreau} writes $h=h^{**}$ as the supremum of its affine minorants, which are affine minorants of $f$; so $h\le f^{**}$. For $f(x)=\min((x-1)^2,(x+1)^2)=(|x|-1)^2$, the envelope is $0$ on $[-1,1]$ and $(|x|-1)^2$ outside:
it is convex, below $f$, and any convex function below $f$ is at most $0$ at $\pm1$, hence at most $0$ on $[-1,1]$, and at most $f$ elsewhere.
\end{solution}

\begin{solution}{ex:slaterfail}
\idea In (a) the multiplier would have to be infinite; in (b) the constraint is degenerate and the Lagrangian can escape along a parabola.
\writeup (a) The only feasible point is $0$, so $p^*=0$, and no point is strictly feasible. $g(\lambda)=\inf_x(x+\lambda x^2)=-\frac1{4\lambda}$ for $\lambda>0$ and $-\infty$ for $\lambda=0$; so
$d^*=\sup_{\lambda>0}(-\frac1{4\lambda})=0=p^*$, not attained. (b) $x_1^2/x_2\le0$ with $x_2>0$ forces $x_1=0$, so $f_0=e^0=1$ and $p^*=1$. For $\lambda\ge0$, along $x_1=t$, $x_2=t^3$,
$e^{-t}+\lambda/t\to0$ as $t\to\infty$, so $g(\lambda)=0$ and $d^*=0$: a duality gap of $1$. Again no point is strictly feasible.
\end{solution}

\begin{solution}{ex:leastnorm}
\idea Stationarity puts $x^*$ in the row space of $A$, and the constraint fixes where.
\writeup Minimise $\frac12\|x\|^2$ subject to $Ax=b$. The KKT conditions are $Ax^*=b$ and $x^*+A^{\mathsf T}\nu^*=0$. Then $-AA^{\mathsf T}\nu^*=b$; $AA^{\mathsf T}$ is invertible because $A$ has full row
rank, so $\nu^*=-(AA^{\mathsf T})^{-1}b$ and $x^*=A^{\mathsf T}(AA^{\mathsf T})^{-1}b$, which is optimal by Theorem~\ref{thm:kkt}. The multiplier is $-(AA^{\mathsf T})^{-1}b$, and $\frac12\|x^*\|^2=\frac12b^{\mathsf T}(AA^{\mathsf T})^{-1}b$
is the dual value. This is the pseudo-inverse solution $A^+b$ of Tome~S1, Theorem~6.7.
\end{solution}

\begin{solution}{ex:waterlevel}
\idea Only the channels below the water line contribute, each with $\ln w$.
\writeup Let $V(P)=\sum_i\ln(\alpha_i+x_i^*)=\sum_{\alpha_i<w}\ln w+\sum_{\alpha_i\ge w}\ln\alpha_i$, with $w=w(P)$ the level. The map $w\mapsto P(w)=\sum_i\max(0,w-\alpha_i)$ is continuous, strictly
increasing for $w>\min\alpha_i$, and piecewise linear with slope $k(w)=\#\{i:\alpha_i<w\}$, so $w(P)$ is continuous, increasing and piecewise linear with slope $1/k$. Between breakpoints,
$V=k\ln w+\text{const}$ and $V'(P)=\frac kw\cdot\frac1k=\frac1w$; $V$ is continuous and its derivative $1/w(P)$ is continuous and decreasing, so $V$ is $C^1$, increasing and concave. The multiplier
$\nu^*=1/w$ of the power constraint is exactly the marginal rate per unit of power.
\end{solution}

\begin{solution}{ex:maxentgauss}
\idea Gibbs' inequality against the Gaussian, whose logarithm is a quadratic, turns the constraint on the variance into a bound on the entropy.
\writeup Let $q(x)=(2\pi\sigma^2)^{-1/2}e^{-x^2/2\sigma^2}$. As in Theorem~\ref{thm:maxent}, $\ln u\le u-1$ gives $\int p\ln\frac qp\le\int(q-p)=0$ over $\{p>0\}$, with equality only if $p=q$ almost
everywhere. Hence
\[
h(p)=-\int p\ln p\le-\int p\ln q=\tfrac12\ln(2\pi\sigma^2)+\frac{1}{2\sigma^2}\int x^2p(x)\dd x=\tfrac12\ln(2\pi\sigma^2)+\tfrac12=\tfrac12\ln(2\pi e\sigma^2),
\]
with equality only for $p=q$. The Gaussian is the maximum entropy law for a given mean and variance; its exponential family form $e^{\theta_1x+\theta_2x^2}$ is that of
Theorem~\ref{thm:maxent} with $\phi(x)=(x,x^2)$.
\end{solution}

\begin{solution}{ex:saddle}
\idea A saddle point is a pair of mutual best responses, and the two values meet at it.
\writeup If $(x^*,y^*)$ is a saddle point, $\min_x\max_y\phi\le\max_y\phi(x^*,y)=\phi(x^*,y^*)=\min_x\phi(x,y^*)\le\max_y\min_x\phi$, and the reverse inequality always holds; so they are
equal, attained at $x^*$ and $y^*$. Conversely, if $x^*$ attains $\min_x\max_y\phi=v$ and $y^*$ attains $\max_y\min_x\phi=v$, then $\phi(x^*,y)\le v$ for all $y$ and $\phi(x,y^*)\ge v$ for all $x$; with
$x=x^*$, $y=y^*$ these give $\phi(x^*,y^*)=v$, and $(x^*,y^*)$ is a saddle point.
\end{solution}

\begin{solution}{ex:rps}
\idea Symmetric games have value $0$; for the finger game, equalise the opponent's options.
\writeup In rock--paper--scissors $M$ is antisymmetric, so $x^{\mathsf T}Mx=0$ and the value is $0$. The uniform strategy gives $M\mathbf1/3=0$, so it guarantees $0$ against everything. If
Rose plays $x\ne$ uniform, some pure answer wins strictly in expectation for Colin (the three expected payments $(M^{\mathsf T}x)_j$ sum to $0$ and are not all zero), so the uniform strategy is
the only optimal one, for both players. In the finger game, with rows and columns indexed by $1$ and $2$ fingers, Rose pays $M=\begin{psmallmatrix}2&-3\\-3&4\end{psmallmatrix}$. If Rose shows one
finger with probability $x$, Colin's columns cost her $2x-3(1-x)=5x-3$ and $-3x+4(1-x)=4-7x$; these are equal at $x=\frac7{12}$, where both are $-\frac1{12}$. By symmetry Colin's
equaliser is also $\frac7{12}$, giving $\min_i(My)_i=-\frac1{12}$. So the value is $-\frac1{12}$: Rose, who looked like the loser, wins $\frac1{12}$ per round on average.
\end{solution}

\section*{Sources and further reading}

\begin{itemize}
\item \emph{A research paper to read.} H.~W.~Kuhn, \emph{``Steiner's'' problem revisited}, in G.~B.~Dantzig and B.~C.~Eaves (eds.), \emph{Studies in Optimization}, MAA Studies in
Mathematics 10 (1974), 52--70. A short and lively paper: read it for the history of the Fermat--Torricelli problem, Fasbender's dual, and the way the general Lagrangian theory
of this chapter recovers both. Then compare with H.~W.~Kuhn, \emph{A note on Fermat's problem}, Mathematical Programming 4 (1973), 98--107, on the convergence of Weiszfeld's
algorithm, and with E.~Weiszfeld, \emph{Sur le point pour lequel la somme des distances de $n$ points donn\'es est minimum}, T\^ohoku Mathematical Journal 43 (1937), 355--386.
\item \emph{The stories.} The epigraph is from R.~T.~Rockafellar, \emph{Lagrange multipliers and optimality}, SIAM Review 35 (1993), 183--238. For Fermat, Torricelli, Cavalieri and
Viviani: Kuhn (1974) above, and B.~S.~Mordukhovich and N.~M.~Nam, \emph{The Fermat--Torricelli problem and Weiszfeld's algorithm in the light of convex analysis}, Journal of
Applied and Numerical Optimization 1 (2019), 205--215. On Karush: H.~W.~Kuhn, \emph{Being in the right place at the right time}, Operations Research 50 (2002), 132--134, and
W.~Karush, \emph{Minima of functions of several variables with inequalities as side conditions}, Master's thesis, University of Chicago, 1939; H.~W.~Kuhn and A.~W.~Tucker,
\emph{Nonlinear programming}, Proceedings of the Second Berkeley Symposium (1951), 481--492. M.~Slater, \emph{Lagrange multipliers revisited}, Cowles Commission Discussion
Paper Math.~403 (1950). E.~T.~Jaynes, \emph{Information theory and statistical mechanics}, Physical Review 106 (1957), 620--630, and \emph{Where do we stand on maximum entropy?},
in \emph{The Maximum Entropy Formalism} (R.~D.~Levine and M.~Tribus, eds.), MIT Press, 1979, for the die with average $4.5$.
\item J.~von Neumann, \emph{Zur Theorie der Gesellschaftsspiele}, Mathematische Annalen 100 (1928), 295--320; \'E.~Borel, \emph{La th\'eorie du jeu et les \'equations int\'egrales \`a noyau
sym\'etrique}, Comptes Rendus 173 (1921), 1304--1308; J.~Ville, \emph{Sur la th\'eorie g\'en\'erale des jeux o\`u intervient l'habilet\'e des joueurs}, in E.~Borel, \emph{Trait\'e du calcul des
probabilit\'es et de ses applications} IV.2 (1938), 105--113; M.~Sion, \emph{On general minimax theorems}, Pacific Journal of Mathematics 8 (1958), 171--176.
\item W.~Fenchel, \emph{On conjugate convex functions}, Canadian Journal of Mathematics 1 (1949), 73--77; J.-J.~Moreau, \emph{Fonctions convexes duales et points proximaux dans un espace
hilbertien}, Comptes Rendus 255 (1962), 2897--2899. T.~M.~Cover and J.~A.~Thomas, \emph{Elements of Information Theory}, 2nd edition, Wiley, 2006, Section~9.4, for water-filling.
\item For the whole chapter: S.~Boyd and L.~Vandenberghe, \emph{Convex Optimization}, Cambridge University Press, 2004, Chapters~2--5 (our proof of Slater's theorem follows
their Section~5.3); R.~T.~Rockafellar, \emph{Convex Analysis}, Princeton University Press, 1970; J.-B.~Hiriart-Urruty and C.~Lemar\'echal, \emph{Fundamentals of Convex Analysis},
Springer, 2001.
\end{itemize}
